\documentclass[12pt]{article}

% --- Packages ---
\usepackage[utf8]{inputenc}
\usepackage[T1]{fontenc}
\usepackage[french]{babel}
\usepackage{amsmath, amssymb}
\usepackage{geometry}
\usepackage{tikz}
\usepackage[most]{tcolorbox}
\usepackage{lmodern}
\usepackage{pgfplots}
\pgfplotsset{compat=1.18}
\usepackage{mdframed}
\usepackage{bm}
\usepackage{xcolor}
\usepackage{sectsty}
\usetikzlibrary{calc}

\usepackage[utf8]{inputenc}
\usepackage[T1]{fontenc}
\usepackage[dvipsnames]{xcolor}
\usepackage{tkz-tab}
\usepackage{multido}
% --- Page ---
\geometry{margin=2cm}

% Footer
\usepackage{fancyhdr}
\pagestyle{fancy}
\cfoot{\thepage}
\rfoot{\textit{G.M.}}
\renewcommand{\headrulewidth}{0pt}
\renewcommand{\footrulewidth}{0pt}

% Pour les normes : 
\newcommand*{\norme}[1]{\left\lVert\overrightarrow{#1}\right\rVert}

% Pour les coordonnées de vecteurs :
\newcommand*{\Coord}[2]{% 
  \ensuremath{
    \begin{pmatrix} 
      #1\\ 
      #2 
    \end{pmatrix}}}

\newcommand*{\Det}[4]{% 
  \ensuremath{
    \begin{vmatrix} 
      #1 & #2 \\
      #3 & #4 
    \end{vmatrix}}}

% % Pour les coordonnées de points :
% \newcommand*{\Coord}[3]{% 
%   \ensuremath{#1\,
%     \begin{pmatrix}
%       #2\\ 
%       #3 
%     \end{pmatrix}}}

\usetikzlibrary{shapes.misc}

\tikzset{cross/.style={cross out, draw=black, minimum size=2*(#1-\pgflinewidth), inner sep=0pt, outer sep=0pt},
%default radius will be 1pt. 
cross/.default={4pt}}

% mettre dans le préambule, après \usepackage{xcolor}
\usepackage{sectsty} % gère facilement la couleur des titres
\setlength{\parindent}{0pt}

\sectionfont{\color{red}\Large\bfseries}     % pour \section*
% \subsectionfont{\color{red}\normalsize\bfseries} % pour \subsection*
\subsectionfont{\color{red}\large\bfseries\hspace{1cm}}

\begin{document}

\begin{center}
    {\LARGE \textbf{Vecteurs dans un repère}}\\[2mm]
    {\large Classe de 2nde}
\end{center}

% ============================
%   PARTIE 1 : REPÈRES
% ============================

\section*{\textcolor{red}{Partie 1 : Repères}}

\vspace{5pt}

\begin{mdframed}[
    linewidth=1pt,
    roundcorner=5pt,
    linecolor=ForestGreen,
    innertopmargin=10pt,
    innerbottommargin=10pt
]
{\color{ForestGreen}\textbf{\underline{Définition :} Repère du plan}}

\vspace{5pt}
Un \textbf{repère du plan} est défini par :
\begin{itemize}
    \item un point $O$, appelé \textbf{origine du repère} ;
    \item deux vecteurs non colinéaires $\overrightarrow{i}$ et $\overrightarrow{j}$.
\end{itemize}
\vspace{5pt}
On note alors le repère :
\[
(O \,;\, \overrightarrow{i}, \overrightarrow{j})
\]
\end{mdframed}

\vspace{10pt}

\underline{\textbf{Les 3 types de repères :}}

\vspace{10pt}

\begin{center}
\begin{tikzpicture}[scale=0.8]

% =======================
% Repère quelconque
% =======================

% Axes (supports de i et j)
\draw[->] (-1,0) -- (4,0) node[right] {$x$};
\draw[->] (-0.33,-1) -- (1.33,4) node[above] {$y$};

% Vecteurs du repère
\draw[->, very thick, blue] (0,0) -- (3,0);
\draw[->, very thick, red] (0,0) -- (1,3);

\node at (0,-0.4) {$O$};
\node[blue] at (1.5,0.4) {$\overrightarrow{i}$};
\node[red] at (0,1.5) {$\overrightarrow{j}$};
\node at (1.2,-1.6) {\textbf{Repère quelconque}};


% =======================
% Repère orthogonal
% =======================

\begin{scope}[xshift=6cm]

% Axes
\draw[->] (-1,0) -- (4,0) node[right] {$x$};
\draw[->] (0,-1) -- (0,4) node[above] {$y$};

% Vecteurs du repère
\draw[->, very thick, blue] (0,0) -- (2.2,0);
\draw[->, very thick, red] (0,0) -- (0,3);

\node at (0,-0.4) {$O$};
\node[blue] at (1,0.4) {$\overrightarrow{i}$};
\node[red] at (-0.3,1.5) {$\overrightarrow{j}$};
\node at (1.2,-1.4) {\textbf{Repère orthogonal}};

\end{scope}

% =======================
% Repère orthonormé
% =======================

\begin{scope}[xshift=12cm]

% Axes
\draw[->] (-1,0) -- (4,0) node[right] {$x$};
\draw[->] (0,-1) -- (0,4) node[above] {$y$};

% Vecteurs du repère
\draw[->, very thick, blue] (0,0) -- (3,0);
\draw[->, very thick, red] (0,0) -- (0,3);

% Carré unité (optionnel mais pédagogique)
\draw[dashed] (3,0) -- (3,3);
\draw[dashed] (0,3) -- (3,3);

\node at (0,-0.4) {$O$};
\node[blue] at (1.5,0.4) {$\overrightarrow{i}$};
\node[red] at (-0.4,1.5) {$\overrightarrow{j}$};
\node at (1.2,-1.4) {\textbf{Repère orthonormé}};

\end{scope}

\end{tikzpicture}
\end{center}


\begin{mdframed}[
    linewidth=1pt,
    roundcorner=5pt,
    linecolor=ForestGreen,
    innertopmargin=10pt,
    innerbottommargin=10pt
]
{\color{ForestGreen}\textbf{\underline{Définitions :} types de repères}}

\vspace{5pt}
\begin{itemize}
    \item Un repère est \textbf{orthogonal} lorsque les vecteurs $\overrightarrow{i}$ et 
    $\overrightarrow{j}$ sont \  \_\_\_\_\_\_\_\_\_\_\_\_\_\_\_.
    \vspace{5pt}
    \item Un repère est \textbf{orthonormé} lorsque les vecteurs $\overrightarrow{i}$ et 
    $\overrightarrow{j}$ sont perpendiculaires \textbf{et de 
    \vspace{10pt}\\même} \_\_\_\_\_\_\_\_\_\_\_\_.
\end{itemize}
\end{mdframed}

\vspace{15pt}

\underline{\textbf{Exemple : coordonnées d’un point}}

\vspace{10pt}

\begin{minipage}{.45\textwidth}
On considère un repère orthonormé $(O \,;\, \overrightarrow{i}, \overrightarrow{j})$.

\vspace{5pt}
Le point $M$ a pour coordonnées $(2\,;3)$.
\end{minipage}
\hfill
\begin{minipage}{.45\textwidth}
\begin{center}
\begin{tikzpicture}[scale=0.6]

% Grille
\draw[step=1cm, gray, very thin] (-1,-1) grid (4,4);

% Axes
\draw[->] (-1,0) -- (4.2,0);
\draw[->] (0,-1) -- (0,4.2);
% Vecteurs du repère
\draw[->, very thick, blue] (0,0) -- (1,0);
\draw[->, very thick, red] (0,0) -- (0,1);

% Point M
\fill (2,3) circle (3pt) node[above right] {$M\ (2\,;3)$};
\node at (-0.4,-0.4) {$O$};
\node[blue] at (1,-0.6) {$\overrightarrow{i}$};
\node[red] at (-0.4,1) {$\overrightarrow{j}$};

\end{tikzpicture}
\end{center}
\end{minipage}

\vspace{10pt}

% ============================
%   PARTIE 2 : Coordonnées d’un vecteur
% ============================

\section*{\textcolor{red}{Partie 2 : Coordonnées d’un vecteur}}

\vspace{5pt}

\underline{\textbf{Exemple 1 :}}

\begin{minipage}[t]{.45\textwidth}%
    \vspace{5pt}
    On considère le repère orthonormé $(O,\overrightarrow{i},\overrightarrow{j})$ ci-contre.
    \vspace{5pt}\\
    On a :
    \[
    \overrightarrow{u} = \_\_\_\_\_\_\_\_\_\_\_\_\_\_
    \]
\end{minipage}%
\hfill
\begin{minipage}[t]{.45\textwidth}%
\vspace{0pt}
\begin{center}
\begin{tikzpicture}[scale=0.8]

\draw[step=1cm, gray, thin] (-1,-1) grid (5,4);
\draw[->] (-1,0) -- (5,0);
\draw[->] (0,-1) -- (0,4);

% Vecteurs du repère
\draw[->, very thick, blue] (0,0) -- (1,0);
\draw[->, very thick, red] (0,0) -- (0,1);
\draw[->, very thick, blue] (1,1) -- (2,1);
\draw[->, very thick, blue] (2,1) -- (3,1);
\draw[->, very thick, blue] (3,1) -- (4,1);
\draw[->, very thick, red] (4,1) -- (4,2);
\draw[->, very thick, red] (4,2) -- (4,3);

\draw[->, very thick, ForestGreen] (1,1) -- (4,3);
\node[ForestGreen, above left] at (3,2) {$\overrightarrow{u}$};
\node at (-0.4,-0.4) {$O$};
\node[blue] at (1,-0.6) {$\overrightarrow{i}$};
\node[red] at (-0.4,1) {$\overrightarrow{j}$};

\end{tikzpicture}
\end{center}
\end{minipage}%

\vspace{10pt}

\underline{Méthode :}
\vspace{5pt}
\begin{itemize}
    \item On part de l’origine du vecteur.
    \item On se déplace de 3 unités sur l'axe des $x$ et de 2 unités sur l'axe des $y$.
    \vspace{15pt}
    \item \textbf{Les coordonnées du vecteur sont donc \_\_\_\_\_\_\_\_\_.}
    \vspace{5pt}
    \item On note généralement les coordonnées d'un vecteur verticalement :
    \Coord{$x$}{$y$}
\end{itemize}

\vspace{20pt}

\underline{\textbf{Exemple 2 :}}

\begin{minipage}[t]{.45\textwidth}%
    \vspace{5pt}
    À nouveau, on se place dans un repère $(O,\overrightarrow{i},\overrightarrow{j})$.
    \vspace{5pt}
    \\On a :
    \[
    \overrightarrow{u} = \_\_\_\_\_\_\_\_\_\_\_\_\_\_
    \]
    Les coordonnées de $\overrightarrow{u}$ sont donc :
    \[
    \Coord{\ \ \ \ \ \ \ }{\ \ \ \ \ \ \ }
    \]
    \begin{mdframed}[
    linewidth=1pt,
    roundcorner=5pt,
    linecolor=Red,
    innertopmargin=10pt,
    innerbottommargin=10pt
    ]
    {\color{Red}\textbf{\underline{À RETENIR :}}}
        Les coordonnées d'un vecteur ne dépendent pas de sa position
        dans le repère, uniquement de sont point de départ et point d'arrivée.
    \end{mdframed}
\end{minipage}%
\hfill
\begin{minipage}[t]{.45\textwidth}%
\begin{center}
\vspace{50pt}
\begin{tikzpicture}[scale=0.8]

\draw[step=1cm, gray, thin] (-1,-2) grid (6,4);
\draw[->] (-1,0) -- (6,0);
\draw[->] (0,-2) -- (0,4);

% Vecteurs du repère
\draw[->, very thick, blue] (0,0) -- (1,0);
\draw[->, very thick, red] (0,0) -- (0,1);
\draw[->, very thick, blue] (1,2) -- (2,2);
\draw[->, very thick, blue] (2,2) -- (3,2);
\draw[->, very thick, blue] (3,2) -- (4,2);
\draw[<-, very thick, red] (4,-1) -- (4,0);
\draw[<-, very thick, red] (4,0) -- (4,1);
\draw[<-, very thick, red] (4,1) -- (4,2);

\draw[->, very thick, ForestGreen] (1,2) -- (4,-1);
\node[ForestGreen, above left] at (3.2,0.5) {$\overrightarrow{u}$};
\node at (-0.4,-0.4) {$O$};
\node[blue] at (1,-0.6) {$\overrightarrow{i}$};
\node[red] at (-0.4,1) {$\overrightarrow{j}$};

\end{tikzpicture}
\end{center}
\end{minipage}%

\vspace{10pt}
\newpage

\begin{mdframed}[
    linewidth=1pt,
    roundcorner=5pt,
    linecolor=ForestGreen,
    innertopmargin=10pt,
    innerbottommargin=10pt
]
{\color{ForestGreen}\textbf{\underline{Définition :}}}
Dans un repère, les \textbf{coordonnées d’un vecteur} $\overrightarrow{u}$ sont les nombres 
$x$ et $y$ tels que $\Coord{x}{y}$ avec :
\[
\overrightarrow{u} = x\overrightarrow{i} + y\overrightarrow{j}
\]
\end{mdframed}

\vspace{10pt}
% =============================================
% ========== EXERCICE D'APPLICATION ============
% =============================================

\begin{mdframed}[
    linewidth=1pt,
    roundcorner=5pt,
    linecolor=RoyalBlue,
    innertopmargin=10pt,
    innerbottommargin=10pt
]
{\color{RoyalBlue}\textbf{\underline{Exercice d'application n°1 :} Lecture graphique de coordonnées.}}
\vspace{10pt}

\begin{minipage}{.45\textwidth}%
\begin{tikzpicture}[scale=0.7]

\draw[step=1cm, gray, thin] (-5,-5) grid (5,5);
\draw[->, thick] (-5,0) -- (5,0);
\draw[->, thick] (0,-5) -- (0,5);

\node at (-0.4,-0.4) {$O$};

% \fill (1,1) circle (2pt) node[below left] {$A$};
% \fill (5,3) circle (2pt) node[above right] {$B$};
% \fill (2,4) circle (2pt) node[above] {$C$};

\end{tikzpicture}
\end{minipage}%
\hfill
\begin{minipage}{.55\textwidth}%
    \begin{itemize}
        \item Dans le repère, placer les points $C\Coord{-1}{-2}$, $D\Coord{-2}{3}$,
        $E\Coord{1}{-4}$ et $F\Coord{4}{-2}$
        \vspace{5pt}
        \item Déterminer les coordonnées des vecteurs 
        $\overrightarrow{CD}$ et $\overrightarrow{EF}$ par lecture graphique.
    \end{itemize}
\end{minipage}%

\end{mdframed}

\vspace{10pt}

\begin{mdframed}[
    linewidth=1pt,
    roundcorner=5pt,
    linecolor=Red,
    innertopmargin=10pt,
    innerbottommargin=10pt
]
{\color{Red}\textbf{\underline{Propriété :}}}
Deux vecteurs sont \textbf{égaux} si et seulement si ils ont les \textbf{mêmes coordonnées}.
\end{mdframed}

\vspace{10pt}

\begin{mdframed}[
    linewidth=1pt,
    roundcorner=5pt,
    linecolor=Red,
    innertopmargin=10pt,
    innerbottommargin=10pt
]
{\color{Red}\textbf{\underline{Propriété :}}}
Soient $A\Coord{x_A}{y_A}$ et $B\Coord{x_B}{y_B}$ deux points. 
\vspace{5pt}
\\Les coordonnées du vecteur $\overrightarrow{AB}$ sont alors :
\[
\overrightarrow{AB}\Coord{x_B - x_A}{y_B - y_A}
\]
\end{mdframed}

\vspace{10pt}

\underline{\textbf{Exemple :}}

\vspace{10pt}
On considère deux points $A\Coord{2}{1}$ et $B\Coord{5}{4}$.

\vspace{5pt}
On a alors : 
\[
\overrightarrow{AB}\Coord{x_B - x_A}{y_B - y_A} \ = \ \overrightarrow{AB} \Coord{5-2}{4-1} \  = \  \overrightarrow{AB}\Coord{3}{3}
\]


\underline{Méthode :}
\vspace{5pt}
\begin{itemize}
    \item On soustrait les coordonnées du point de départ à celles du point d’arrivée.
\end{itemize}


\newpage
% =============================================
% ========== EXERCICE D'APPLICATION ============
% =============================================

\begin{mdframed}[
    linewidth=1pt,
    roundcorner=5pt,
    linecolor=RoyalBlue,
    innertopmargin=10pt,
    innerbottommargin=10pt
]
{\color{RoyalBlue}\textbf{\underline{Exercice d'application n°2 :} Déterminer les coordonnées d'un 
vecteur par le calcul.}}

\vspace{5pt}

Soient les points : 

$A\Coord{2}{1}$, $B\Coord{5}{3}$, $C\Coord{-1}{-2}$, $D\Coord{-2}{3}$, $E\Coord{1}{-4}$ et $F\Coord{4}{-2}$.
\vspace{5pt}

Calculer les coordonnées des vecteurs $\overrightarrow{AB}$, $\overrightarrow{CD}$ 
et $\overrightarrow{EF}$.

\begin{mdframed}[
    linewidth=1pt,
    roundcorner=5pt,
    linecolor=black,
    innertopmargin=10pt,
    innerbottommargin=10pt,
    backgroundcolor=gray!15
]
\underline{\textbf{Méthode :}}
\begin{itemize}
    \item appliquer la propriété sur les coordonnées de vecteurs
\end{itemize}
\end{mdframed}

\end{mdframed}

\vspace{10pt}

\begin{mdframed}[
    linewidth=1pt,
    roundcorner=5pt,
    linecolor=Red,
    innertopmargin=10pt,
    innerbottommargin=10pt
]
{\color{Red}\textbf{\underline{Propriété :}}}

Soient deux vecteurs $\overrightarrow{u}\Coord{x}{y}$ et $\overrightarrow{v}\Coord{x'}{y'}$ et un réel 
$k$.

On a :
\begin{itemize} 
\renewcommand{\labelitemi}{$\bullet$}
    \item $\overrightarrow{u} + \overrightarrow{v} \Coord{x+x'}{y+y'}$ \hspace{2cm} $\bullet$
    $\overrightarrow{u} - \overrightarrow{v} \Coord{x-x'}{y-y'}$
    \vspace{10pt}
    \item $k \overrightarrow{u} \Coord{kx}{ky}$ \hspace{3.3cm} $\bullet$ $-\overrightarrow{u}\Coord{-x}{-y}$
\end{itemize}

\end{mdframed}

% =============================================
% ========== EXERCICE D'APPLICATION ============
% =============================================

\begin{mdframed}[
    linewidth=1pt,
    roundcorner=5pt,
    linecolor=RoyalBlue,
    innertopmargin=10pt,
    innerbottommargin=10pt
]
{\color{RoyalBlue}\textbf{\underline{Exercice d'application n°3 :}
Appliquer les formules sur les coordonnées de vecteurs.}}
\vspace{10pt}

Soient les points : 

\vspace{5pt}

$A\Coord{2}{1}$, $B\Coord{5}{3}$, $C\Coord{-1}{-2}$ et $D\Coord{-2}{3}$

\vspace{10pt}

Calculer les coordonnées des vecteurs :
\vspace{5pt}

\hspace{1cm} $\bullet \  3 \overrightarrow{AB}$ \hspace{2cm} $\bullet \ 4 \overrightarrow{CD}$ \hspace{2cm} 
$\bullet \ 3 \overrightarrow{AB} - 4 \overrightarrow{CD}$

\begin{mdframed}[
    linewidth=1pt,
    roundcorner=5pt,
    linecolor=black,
    innertopmargin=10pt,
    innerbottommargin=10pt,
    backgroundcolor=gray!15
]
\underline{\textbf{Méthode :}}
\begin{itemize}
    \item calculer d'abord le vecteur
    $\overrightarrow{AB} \Coord{x_B-x_A}{y_B-y_A}$
    \item appliquer la propriété pour calculer le vecteur final 
    $3 \overrightarrow{AB}$
    \item faire de même avec les autres calculs de vecteurs
\end{itemize}
\end{mdframed}

\vspace{1cm}
\newpage
{\color{RoyalBlue}\textbf{\underline{Exercice d'application n°4 :}
Calculer les coordonnées d'un point défini par une égalité 
vectorielle.}}
\vspace{10pt}

Soient les points : 

\vspace{5pt}

$A\Coord{1}{2}$, $B\Coord{-4}{3}$ et $C\Coord{1}{-2}$.

\vspace{10pt}

Déterminer les coordonnées du point $D$ tel que $ABCD$ soit un 
parallélogramme.
\vspace{5pt}

\begin{mdframed}[
    linewidth=1pt,
    roundcorner=5pt,
    linecolor=black,
    innertopmargin=10pt,
    innerbottommargin=10pt,
    backgroundcolor=gray!15
]
\underline{\textbf{Méthode :}}
\begin{itemize}
    \item $ABCD$ est un parallélogramme si et seulement si 
    \_\_\_\_\_\_\_\_\_\_\_\_\_\_\_\_\_\_.
    \vspace{5pt}
    \item en posant $\Coord{x_D}{y_D}$ les coordonnées du point $D$, on peut calculer 
    les coordonnées de ces 2 vecteurs.
    \vspace{5pt}
    \item on écrit alors \_\_\_\_\_\_\_\_\_\_\_\_\_\_\_\_\_\_\_\_\_\_\_\_\_\_\_\_\_\_\_\_\_\_\_\_\_\_\_\_.
    
    \vspace{5pt}
    Soit $1 - x_D = -5$ et $-2 -y_D = 1$.
    \vspace{5pt}
    \item on résout alors pour trouver les coordonnées de $D$.
\end{itemize}
\end{mdframed}

\end{mdframed}

\vspace{20pt}

% ============================
%   PARTIE 3 : Colinéarité de deux vecteurs
% ============================

\section*{\textcolor{red}{Partie 3 : Colinéarité de deux vecteurs}}

\subsection*{\textcolor{red}{1. Critère de colinéarité}}

\begin{mdframed}[
    linewidth=1pt,
    roundcorner=5pt,
    linecolor=Red,
    innertopmargin=5pt,
    innerbottommargin=10pt
]
{\color{Red}\textbf{\underline{Propriété :}}}
Soit deux vecteurs $\overrightarrow{u}\Coord{x}{y}$ et $\overrightarrow{v}\Coord{x'}{y'}$.
\\Dire que $\overrightarrow{u}$ et $\overrightarrow{v}$ sont colinéaires revient à dire 
que :
\[
xy'-yx' = 0
\]
\end{mdframed}

% \vspace{20pt}

\underline{Remarque :} Dire que $\overrightarrow{u}$ et $\overrightarrow{v}$ sont colinéaires 
revient à dire que les coordonnées des deux vecteurs sont proportionnelles soit :
$xy' = yx'$.

\begin{mdframed}[
    linewidth=1pt,
    roundcorner=5pt,
    linecolor=RoyalBlue,
    innertopmargin=10pt,
    innerbottommargin=10pt
]
{\color{RoyalBlue}\textbf{\underline{Démonstration au programme :}}}
\vspace{10pt}

Supposons $\overrightarrow{u}$ et $\overrightarrow{v}$ non nuls (égalité évidente 
si l'un des deux est nul).
\vspace{5pt}

\underline{1\up{er} sens :} Si $\overrightarrow{u}$ et $\overrightarrow{v}$ colinéaires 
alors $xy'-yx' = 0$.

% \vspace{5pt}

Si $\overrightarrow{u}$ et $\overrightarrow{v}$ colinéaires alors il existe un réel 
$k$ tel que $\overrightarrow{u} = k \overrightarrow{v}$.

Les coordonnées des vecteurs sont donc proportionnelles et le tableau ci-dessous est 
un tableau de proportionnalité :
\begin{center}
    \begin{tabular}{|c|c|}
        \hline
        x & x' \\
        \hline
        y & y' \\
        \hline
    \end{tabular}
\end{center}

Donc (produit en croix) : $xy' = yx'$ ou encore $xy' - yx' = 0$.

\vspace{5pt}
\underline{2\up{ème} sens :} Si $xy'-yx' = 0$  
alors $\overrightarrow{u}$ et $\overrightarrow{v}$ sont colinéaires.

\vspace{5pt}

Le vecteur $\overrightarrow{v}$ étant non nul alors au moins l'une de ses coordonnées 
est non nulle. Supposons $x'$ non nul.

On pose alors $k = \dfrac{x}{x'}$. L'égalité $xy'-yx' = 0$ devient :
\begin{align*}
yx' &= xy' \\
y &= \dfrac{xy'}{x'} \\
y &= ky'
\end{align*}

Comme on a déjà $x = kx'$, on en déduit que $\overrightarrow{u} = k \overrightarrow{v}$.


\end{mdframed}

\vspace{10pt}


\begin{mdframed}[
    linewidth=1pt,
    roundcorner=5pt,
    linecolor=RoyalBlue,
    innertopmargin=10pt,
    innerbottommargin=10pt
]
{\color{RoyalBlue}\textbf{\underline{Exercice d'application n°5 :} 
Vérifier si deux vecteurs sont colinéaires.}}
\vspace{10pt}

Dans chaque cas, vérifier si les vecteurs $\overrightarrow{u}$ et $\overrightarrow{v}$ 
sont colinéaires.

a) $\overrightarrow{u}\Coord{4}{-7}$ et $\overrightarrow{v}\Coord{-12}{21}$
\hspace{1cm} b) $\overrightarrow{u}\Coord{5}{-2}$ et $\overrightarrow{v}\Coord{15}{-7}$

\begin{mdframed}[
    linewidth=1pt,
    roundcorner=5pt,
    linecolor=black,
    innertopmargin=10pt,
    innerbottommargin=10pt,
    backgroundcolor=gray!15
]
\underline{\textbf{Méthode :}}
\begin{itemize}
    \item appliquer directement la propriété aux coordonnées des vecteurs
    \item on peut parfois remarquer directement une relation de proportionnalité 
    entre les vecteurs
    \item avec la propriété, on trouve :\\a) les vecteurs sont bien colinéaires (on 
    remarque d'ailleurs  
    $\overrightarrow{v} = $\_\_\_\_\_\_\_\_\_\_)
    \\b) les vecteurs ne sont pas colinéaires.
\end{itemize}
\end{mdframed}

\end{mdframed}

\vspace{20pt}

\subsection*{\textcolor{red}{2. Déterminant de deux vecteurs}}

\vspace{5pt}

\begin{mdframed}[
    linewidth=1pt,
    roundcorner=5pt,
    linecolor=ForestGreen,
    innertopmargin=10pt,
    innerbottommargin=10pt
]
{\color{ForestGreen}\textbf{\underline{Définition :} Déterminant de deux vecteurs}}

\vspace{5pt}
Soit $\overrightarrow{u}\Coord{x}{y}$ et $\overrightarrow{v}\Coord{x'}{y'}$.

\vspace{5pt}
Le nombre $xy' - yx'$ est appelé \textbf{déterminant} des vecteurs $\overrightarrow{u}$ 
et $\overrightarrow{v}$.

\vspace{5pt}
On note :
\[
det(\overrightarrow{u} ; \overrightarrow{v}) = \Det{x}{x'}{y}{y'} = xy' - yx'
\]
\end{mdframed}

\vspace{10pt}

\begin{mdframed}[
    linewidth=1pt,
    roundcorner=5pt,
    linecolor=Red,
    innertopmargin=10pt,
    innerbottommargin=10pt
]
{\color{Red}\textbf{\underline{Propriété :}}}
Dire que $\overrightarrow{u}$ et $\overrightarrow{v}$ sont colinéaires revient à dire que 
$det(\overrightarrow{u};\overrightarrow{v}) = 0$.
\end{mdframed}

\vspace{10pt}

\begin{mdframed}[
    linewidth=1pt,
    roundcorner=5pt,
    linecolor=RoyalBlue,
    innertopmargin=10pt,
    innerbottommargin=10pt
]
{\color{RoyalBlue}\textbf{\underline{Exercice d'application n°6 :} 
Vérifier si deux vecteurs sont colinéaires à l'aide du déterminant.}}
\vspace{10pt}

Dans chaque cas, vérifier si les vecteurs $\overrightarrow{u}$ et $\overrightarrow{v}$ 
sont colinéaires.

a) $\overrightarrow{u}\Coord{-6}{10}$ et $\overrightarrow{v}\Coord{9}{-15}$
\hspace{1cm} b) $\overrightarrow{u}\Coord{4}{9}$ et $\overrightarrow{v}\Coord{11}{23}$

\begin{mdframed}[
    linewidth=1pt,
    roundcorner=5pt,
    linecolor=black,
    innertopmargin=10pt,
    innerbottommargin=10pt,
    backgroundcolor=gray!15
]
\underline{\textbf{Méthode :}}
\begin{itemize}
    \item on écrit le déterminant en mettant les coordonnées aux bons endroits
    \item on résout le calcul du déterminant
    \item on conclut
    \item \underline{Résultats (comparer après calcul) :}\\
    a) $\overrightarrow{u}$ et $\overrightarrow{v}$ sont colinéaires.
    \\b) $\overrightarrow{u}$ et $\overrightarrow{v}$ ne sont pas colinéaires.
\end{itemize}
\end{mdframed}

\end{mdframed}

\vspace{20pt}

\subsection*{\textcolor{red}{3. Applications à la géométrie}}

\vspace{5pt}

\begin{mdframed}[
    linewidth=1pt,
    roundcorner=5pt,
    linecolor=Red,
    innertopmargin=10pt,
    innerbottommargin=10pt
]
{\color{Red}\textbf{\underline{Propriété :}}}

\vspace{5pt}
1) Dire que les droites $(AB)$ et $(CD)$ sont \textbf{parallèles} revient à dire que les 
vecteurs $\overrightarrow{AB}$ et $\overrightarrow{CD}$ sont colinéaires.

\vspace{5pt}
2) Dire que les points $A$, $B$ et $C$ sont \textbf{alignés} revient à dire que les vecteurs 
$\overrightarrow{AB}$ et $\overrightarrow{AC}$ sont colinéaires.
\end{mdframed}

\vspace{5pt}

\begin{mdframed}[
    linewidth=1pt,
    roundcorner=5pt,
    linecolor=RoyalBlue,
    innertopmargin=10pt,
    innerbottommargin=10pt
]
{\color{RoyalBlue}\textbf{\underline{Exercice d'application n°7 :} 
Appliquer la colinéarité à des question de géométrie.}}
\vspace{10pt}

On considère les points $A\Coord{-1}{1}$, $B\Coord{3}{2}$, $C\Coord{-2}{-3}$, 
$D\Coord{6}{-1}$ et $E\Coord{5}{0}$.

\vspace{5pt}
a) Démontrer que les droites $(AB)$ et $(CD)$ sont parallèles.
\\b) Démontrer que les points $E$, $B$ et $D$ sont alignés.

\begin{mdframed}[
    linewidth=1pt,
    roundcorner=5pt,
    linecolor=black,
    innertopmargin=10pt,
    innerbottommargin=10pt,
    backgroundcolor=gray!15
]
\underline{\textbf{Méthode :}}
\vspace{5pt}
\\a) 
\begin{itemize}
    \item on utilise la propriété : il faut prouver que \_\_\_\_\_\_ et \_\_\_\_\_\_ 
    sont colinéaires (ou non)

    \vspace{10pt}
    \item pour cela on utilise le calcul du \_\_\_\_\_\_\_\_\_\_\_\_\_\_\_
    
    \vspace{10pt}
    \item on peut remarquer aussi que les coordonnées de $\overrightarrow{AB}$ et $\overrightarrow{CD}$ 
    sont proportionnelles et donc les vecteurs sont colinéaires.
    
    \vspace{10pt}
    On a en effet : $\overrightarrow{CD} =$ \_\_\_\_\_\_\_\_\_\_\_\_\_\_\_
\end{itemize}

\vspace{5pt}
b) 
\begin{itemize}
    \item idem on utilise la propriété
\end{itemize}
\end{mdframed}

\end{mdframed}


\vspace{10pt}

% ============================
%   PARTIE 4 : Coordonnées du milieu d'un segment
% ============================

\section*{\textcolor{red}{Partie 4 : Coordonnées du milieu d'un segment}}

\vspace{5pt}

\begin{mdframed}[
    linewidth=1pt,
    roundcorner=5pt,
    linecolor=Red,
    innertopmargin=10pt,
    innerbottommargin=10pt
]
{\color{Red}\textbf{\underline{Propriété :}}}

\vspace{5pt}
Soit $A\Coord{x_A}{y_A}$ et $B\Coord{x_B}{y_B}$.
\\Le milieu $M$ du segment $\left[AB\right]$ a pour 
coordonnées : $\renewcommand*{\arraystretch}{2}\begin{pmatrix} 
      \dfrac{x_A+x_B}{2}\\
      \dfrac{y_A+y_B}{2}
    \end{pmatrix}$
\end{mdframed}

\vspace{5pt}

\begin{mdframed}[
    linewidth=1pt,
    roundcorner=5pt,
    linecolor=RoyalBlue,
    innertopmargin=10pt,
    innerbottommargin=10pt
]
{\color{RoyalBlue}\textbf{\underline{Exercice d'application n°8 :} 
Calculer les coordonnées d'un milieu.}}
\vspace{5pt}

On considère les points $A\Coord{2}{3}$, $B\Coord{-2}{1}$ et $C\Coord{3}{-1}$.

\vspace{5pt}
Calculer les coordonnées de $M$, $N$ et $P$ milieux respectifs de 
$\left[AB\right]$, $\left[AC\right]$ et $\left[BC\right]$.

\begin{mdframed}[
    linewidth=1pt,
    roundcorner=5pt,
    linecolor=black,
    innertopmargin=10pt,
    innerbottommargin=10pt,
    backgroundcolor=gray!15
]
\underline{\textbf{Méthode :}}
\begin{itemize}
    \item on applique la propriété aux coordonnées des points
\end{itemize}
\end{mdframed}

\end{mdframed}


\vspace{10pt}
\newpage
% ============================
%   PARTIE 5 : Distance dans un repère orthonormé
% ============================

\section*{\textcolor{red}{Partie 5 : Distance dans un repère orthonormé}}

\vspace{5pt}

\begin{mdframed}[
    linewidth=1pt,
    roundcorner=5pt,
    linecolor=Red,
    innertopmargin=10pt,
    innerbottommargin=10pt
]
{\color{Red}\textbf{\underline{Propriété :}}}

\vspace{5pt}
Soit deux points $A\Coord{x_A}{y_A}$ et $B\Coord{x_B}{y_B}$ dans un 
repère \textbf{orthonormé}.

\vspace{5pt}
La distance $AB$ (ou la norme $\norme{AB}$) est :
\[
AB = \norme{AB} = \sqrt{(x_B - x_A)^2 + (y_B-y_A)^2}
\]
\end{mdframed}

\underline{Remarque :} Cette propriété est une conséquence 
du théorème de Pythagore.

\vspace{5pt}

\begin{mdframed}[
    linewidth=1pt,
    roundcorner=5pt,
    linecolor=RoyalBlue,
    innertopmargin=10pt,
    innerbottommargin=10pt
]
{\color{RoyalBlue}\textbf{\underline{Exercice d'application n°9 :} 
Calculer.}}
\vspace{5pt}

On considère les points $A\Coord{3}{2}$ et $B\Coord{2}{-2}$ dans un repère orthonormé.

\vspace{5pt}
Calculer la distance $AB$.

\begin{mdframed}[
    linewidth=1pt,
    roundcorner=5pt,
    linecolor=black,
    innertopmargin=10pt,
    innerbottommargin=10pt,
    backgroundcolor=gray!15
]
\underline{\textbf{Méthode :}}
\begin{itemize}
    \item on applique la propriété aux coordonnées des points
\end{itemize}
\end{mdframed}

\end{mdframed}


\end{document}
